\documentclass[11pt,a4paper]{article}
\usepackage[utf8]{inputenc}
\usepackage[T1]{fontenc}
\usepackage{lmodern}
\usepackage[margin=2.5cm]{geometry}
\usepackage{amsmath,amssymb}
\usepackage{booktabs}
\usepackage{longtable}
\usepackage{enumitem}
\usepackage{hyperref}
\usepackage{xcolor}
\usepackage{titlesec}
\usepackage{framed}

\hypersetup{
    colorlinks=true,
    linkcolor=blue!60!black,
    citecolor=blue!60!black,
    urlcolor=blue!60!black
}

\newcommand{\keep}{\textsc{retained}}
\newcommand{\reform}{\textsc{reformulated}}
\newcommand{\retract}{\textcolor{red!65!black}{\textsc{retracted}}}
\newcommand{\dormant}{\textsc{dormant}}
\newcommand{\executed}{\textcolor{green!35!black}{\textsc{executed}}}

\title{\textbf{The Pharmacist's Cipher\,II}\\[0.4em]
\large Status Report and Partial Retraction\\[0.3em]
\normalsize Version 2.0}

\author{Alessandro Placa\thanks{Corresponding author: \texttt{alessandro.placa@uroboro.io}}\\
\textit{Uroboro.io}}

\date{July 2026\\[0.3em]
\small Supersedes: Placa (2026), ``The Pharmacist's Cipher\,II'', v1.2,\\
\small DOI: 10.5281/zenodo.19228502}

\begin{document}
\maketitle

\begin{framed}
\noindent\textbf{What this document is.} In March 2026 I published a companion report to
``The Pharmacist's Cipher'' (DOI: 10.5281/zenodo.19197846) setting out a working morphological
model for MS\,408, a set of visual-textual convergences, and 39 dated, falsifiable predictions.

Four months of further work have falsified or superseded substantial parts of that report. This
document is not a revision of the companion. It is a section-by-section status assessment, to be
read \emph{before} the original, which is included unchanged as the second file in this record.

I am retracting the therapeutic-class system, the botanical identifications, the chromatic code,
and the cardinal reading of folio f85r. The structural and distributional findings stand, in some
cases in reformulated terms. The prediction register stands as a dated record, and one of its
entries has since been executed with results that could not have been obtained without it.

Nothing here was prompted by external criticism. The falsifications are my own, produced by the
tests the original report asked for.
\end{framed}

\section{Why this document exists rather than a revised edition}

The companion report was explicitly built as an atlas of disciplined hypotheses, with an epistemic
label on every section and a declared falsification criterion for each claim. That design has now
done its work: enough of those criteria have been met that editing the text in place would produce
a document whose surviving claims could no longer be distinguished from its retracted ones.

A reader arriving at the original today cannot tell which of its 29 pages are still load-bearing.
That is the problem this status report solves, and it solves it by leaving the original intact and
supplying the key to it, rather than by silently rewriting history. The original remains citable at
its own DOI, and any claim of intellectual priority I make rests on its March 2026 date, not on a
later edit.

\section{Section-by-section status}

Five verdicts are used. \keep{}: the finding stands as stated. \reform{}: the measurement stands
but the interpretation attached to it has been replaced. \retract{}: withdrawn, do not use.
\dormant{}: registered but never tested, and with no current supporting evidence, so it should be
read as an open conjecture rather than a result. \executed{}: a declared research requirement that
has since been carried out.

\begin{longtable}{@{}p{4.9cm}p{2.4cm}p{7.6cm}@{}}
\toprule
\textbf{Section} & \textbf{Status} & \textbf{Assessment} \\
\midrule
\endfirsthead
\toprule
\textbf{Section} & \textbf{Status} & \textbf{Assessment} \\
\midrule
\endhead
\bottomrule
\endfoot

2.1.1 Prefix families and distributional profiles & \reform{} &
The distributional facts hold, including the anomalous $o$-prefix rate for $s$- (4.5\% against
54--71\%). The framing as ``four solvent-prefix families'' is now a reading, not a fact: the
substance glosses are held at structural-solid / open-reading status. The $s$- anomaly is better
accounted for by the finding that $s$ and $d$- do not host the gallows morpheme at all. \\
\addlinespace

2.1.2 Segmentation claims & \reform{} &
Mixed. The decomposition $dy = d+y$ holds. The unitary status of $ch$ and $sh$ holds and has been
strengthened: $ckh$, $cth$, $cph$, $cfh$ are now treated as atomic. The coalescence rules
$qo + ok \rightarrow qok$ hold and $qo$ has been independently validated across 5{,}181
occurrences. The falsification of the calligraphic model for simple versus ornate gallows holds,
but its replacement has changed (see 2.1.3). The reading of \textit{-am} as a maceration operator
is withdrawn: \textit{-am} marks completion of a step. \\
\addlinespace

2.1.3 Thermal asymmetry across gallows classes & \reform{} &
The differential co-occurrence of $ch$ and $sh$ with $k/t$ versus $f/p$ is a real measurement and
was an early detection of a genuine split. The reading was wrong. The axis is not thermal but a
two-by-two factorization of process domain (wet versus dry), crossed with three graphic grades
carried in the feet of the glyph. A further asymmetry has since been measured on the same split:
$k$ and $t$ behave lexically, $p$ and $f$ behave positionally. \\
\addlinespace

2.1.4 Distributional complementarity: ornate $c\_h$ and bench characters & \keep{} &
The complementarity table stands, and the paradigmatic correspondence between ornate forms and
simple $ch$/$sh$ forms was the first evidence for what is now the atomic treatment of the ornate
bench characters. The specific interpretation as compression of a \emph{thermal} marker falls with
2.1.3; the compression finding itself does not. \\
\addlinespace

2.1.5 The 224/224 coverage result & \reform{} &
Superseded by stronger measures rather than falsified. Coverage of first tokens was always offered
as compatibility evidence, and it is a weak instrument: a sufficiently flexible agglutinative
template will cover almost anything. It has been replaced by direct measurements of the shared
skeleton, of the closed prefix inventory, and of the one-gallows-per-word constraint at 97.7\%. \\
\addlinespace

2.2.1 Solvent-prefix glosses & \reform{} &
Demoted from working glosses to open readings. Only $l$ as spirit or alcohol retains independent
support. The others are compatible with the distributional evidence but are not entailed by it,
and at least three non-pharmaceutical frames survive the same evidence. \\
\addlinespace

2.2.2 Therapeutic class glosses (P-CALM, K-EXCIT, T-PROT, F-ABORT) & \retract{} &
Withdrawn in full. The four gallows identities do not encode therapeutic classes. They encode
process domains in a two-by-two-by-three system, and the third dimension of that system is not
visible in the transcription alphabet at all. Every claim in this report that uses the four class
labels as a column falls with them. \\
\addlinespace

2.2.3 Morpheme glosses & \reform{} &
Mixed, item by item. Retained and strengthened: $o$ as a locative morpheme, $d$ as a contextual
referent, $l$ as spirit, \textit{-y} as an action marker. Withdrawn: $a$ as a surface preposition
(it marks extraction), \textit{pol} as laudanum, \textit{-am} as maceration. Reformulated: the
$e/ee/eee$ series is a count of processing cycles rather than a quantity gradient, and those cycles
have since been shown to attach to the concentrated grade, which kills the reading of cycles as
distillation. Dormant: $r$ as a purity marker, the $eo$ system, \textit{raiin}. \\
\addlinespace

2.3 Two orthogonal pairs: $sh/ch$ and $s/l$ & \reform{} &
The claim that the system carries two independent axes survives and has been generalized. The
specific assignment of those axes (thermal versus solvent) does not, and $s$ has since been shown
to behave unlike the other three members. \\
\addlinespace

2.4 The verbal paradigms $qok$-, $qot$-, $qol$- & \keep{} &
Stands, and is now better supported than when written. \\
\addlinespace

3.1 f86v4 structural anomaly & \keep{} &
The structural facts (token density, hapax ratio, bipartite layout, line-1 uniqueness) are
correct and remain interesting. \\
\addlinespace

3.2 f86v4 operational reading (North, East, South, West cycle) & \retract{} &
Withdrawn. The folio is better described as a hybrid of code and data with a metadocumentary
function; its rings contain no tokens of the family the cardinal reading requires. No sequential
production cycle is sustainable on current evidence. \\
\addlinespace

3.3 Exploratory interpretations and implied social context & \dormant{} &
Registered in March as explicitly speculative and never tested since. The mnemonic-song reading and
the inferences drawn from it should be treated as conjecture with no supporting work behind them.
No biographical or identitarian claim was made then and none is made now. \\
\addlinespace

4 The chromatic code & \retract{} &
Withdrawn. It depends on the botanical identifications retracted below, it was never replicated
under a pre-registered protocol, and the direction of inference it relies on cannot be audited
after the fact. \\
\addlinespace

5.1 The 3x3 rosette grid: cardinal solvent domains & \retract{} &
Withdrawn, and this one is a genuine falsification rather than a withdrawal for weakness. The
report assigned the west rosettes to the acetous domain while noting that no textual confirmation
existed. The west band has since been measured, on two independent transcribers with a
pre-declared direction, and it belongs to a different family entirely (Fisher odds ratio 4.78,
$p = 0.0004$). The east-west axis runs between two families, neither of which is the one predicted.
\\
\addlinespace

5.2 Angular rosettes as administration ports & \retract{} &
Withdrawn. The measured organization of the board assigns different roles to those positions. \\
\addlinespace

5.3 Dosage bundles (\textit{fasci}) & \dormant{} &
Never tested. The retrovalidation offered in support was a single correlation on a small
partition and should not be relied on. \\
\addlinespace

5.4 Implicit base vectors & \retract{} &
Withdrawn: it is stated in terms of the therapeutic classes retracted at 2.2.2. \\
\addlinespace

6.1 The sealed-envelope protocol & \dormant{} &
The protocol was designed and applied by its author and is therefore not a blind test, as the
report itself stated. It produced no result that survived. It is left on record as a method
proposal, not as a validated instrument. \\
\addlinespace

6.2--6.4 Botanical identifications (17 plants) & \retract{} &
Withdrawn in full, including the high-convergence table, the Solanaceae quartet, and the
lower-confidence list. The identifications inherit the retracted therapeutic classes as their
organizing column and the retracted chromatic code as their confirmation channel. No plant
identification from this report should be cited. \\
\addlinespace

6.5 The Bax retrodiction test & \retract{} &
Withdrawn. It scores predictions that are themselves retracted, against a rubric written by the
author, on a sample of eight. The report already flagged these limits; they are now decisive. \\
\addlinespace

11 The falsification register (26 hypotheses) & \keep{} &
Stands, and is extended by this document. \\
\addlinespace

12 Epistemological framework & \keep{} &
Stands. The self-assessment it contains, 95 out of 100 for internal coherence against 63 out of 100
for external probative force, was accurate, and the events recorded here are what that gap looks
like when it is honoured. \\
\addlinespace

14 Cautious notes on the implied producer & \dormant{} &
Untested inference. Should not be cited. \\
\end{longtable}

\section{Status of the 39 registered predictions}

The prediction register is the part of the companion report that was designed to age, and it has
aged as intended: each entry carried a falsification criterion, and the criteria did their work.

\begin{center}
\begin{tabular}{@{}lcl@{}}
\toprule
\textbf{Status} & \textbf{Count} & \textbf{Entries} \\
\midrule
\keep{}     & 7  & P3, P4, P5, P7, P10, P12, P13 \\
\reform{}   & 5  & P1, P2, P6, P9, P30 \\
\retract{}  & 13 & P8, P11, P14, P15, P16, P18, P19, P20, P21, P29, P31, P32, P33 \\
\dormant{}  & 12 & P17, P22, P24, P25, P26, P27, P28, P34, P35, P36, P37, P39 \\
\executed{} & 2  & P23 (in part), P38 \\
\midrule
\textbf{Total} & \textbf{39} & \\
\bottomrule
\end{tabular}
\end{center}

\subsection*{P38, and why the register was worth keeping}

Prediction P38 was not a claim about the manuscript. It was a statement about what the research
lacked. It said that a paleographic catalogue of the gallows variants, encoding each as a directed
stroke sequence, was a necessary precondition for testing the framework, and that it did not exist.

That catalogue now exists. Twelve canonical variants were identified, anchored in the feet of the
glyph, and 6{,}143 gallows have been traced at ductus resolution across the manuscript with a
public, versioned dataset behind them.

What it produced was not a confirmation of the March framework. It was a finding the March
framework could not have reached and did not anticipate: the variant grade behaves as a third
information channel, carried at the level of the individual occurrence, largely independent of the
word it sits in and not predictable from its context. A substantial share of the information in the
gallows slot is therefore invisible to any transcription alphabet that collapses the variants,
including the one this report used throughout.

This is the reason the companion report should not be retracted in full. Its interpretive layer was
wrong in the places listed above. Its methodological layer identified, in March, the specific
instrument whose absence was blocking progress, and building that instrument is what produced the
result that now supersedes much of the report.

\section{What this does not change}

\begin{itemize}[leftmargin=1.4em]
\item \textbf{Priority.} The predictions in the original report are dated March 2026 and remain so.
Retraction of some of them does not alter the date of the others.
\item \textbf{Paper I.} The six statistical tests are a separate document and are assessed
separately. Four of them stand as measurements. One stands with a rewritten interpretation. One,
the morpheme-frequency volume hierarchy, is being withdrawn: it depends on a fifth family that the
current model does not recognize as belonging to the same series, and without it the rank
correlation does not survive. A corrected edition is in preparation.
\item \textbf{The open problem.} MS\,408 is not deciphered, and nothing in the original report or
in this one claims otherwise.
\end{itemize}

\section{Where the work goes}

A third paper is in preparation, deliberately agnostic about interpretation. It treats only the
sign system and the mathematics of the corpus: the twelve ductus variants and their factorization,
the non-random signal they carry, the complete positioned transliteration of the Nine-Rosettes
foldout, and the information architecture of the text measured against a genre-matched and
size-matched real-language control. It will make no pharmaceutical claim of any kind. It is
scheduled to follow the completion of the full positioned transcription rather than to precede it,
because the transcription is not supporting documentation for that paper: it is the instrument that
makes the third information channel measurable at all.

Readers who want to check the present work rather than the retracted work should wait for that
paper, or consult the public datasets in the meantime.

\section*{Files in this record}

\begin{enumerate}[leftmargin=1.6em]
\item This status report (v2.0).
\item \texttt{pharmacists-cipher-II\_v1.2\_superseded.pdf}: the original companion report of March
2026, unaltered, preserved for citation and for the record. Read it through the table in Section 2.
\end{enumerate}

\section*{Acknowledgments}

Computational analysis was performed with the assistance of Claude (Anthropic). All interpretive
decisions, and all errors including the ones retracted here, are the author's responsibility.

\end{document}
